\documentclass[instructions.tex]{subfiles}
\begin{document}
 
\chapter{Du zèle pour le salut des enfans.}
 
\primo Quelques mots de catéchisme qu’on apprend dans sa jeunesse, ne sont pas suffisans pour l’instruction d’un chrétien. Si tant de gens vivent dans le désordre, c’est parce qu’ils n’ont jamais été instruits à fond des maximes de la religion dans leur jeunesse, et qu’on a laissé croître leurs vices avec l’âge. L’attachement au monde et l’esprit d’intérêt aveuglent tellement les parens, qu’ils ont bien plus de zèle pour établir leurs enfans que pour les sanctifier.
 
Vous faites tort à vos enfans, si vous les faites riches sans les faire bons chrétiens; il vaut mieux qu’ils manquent de biens, que de manquer de vertus et d’instruction. Les biens, sans la crainte de Dieu, ne serviront qu’à les rendre plus vicieux, et à vous perdre avec eux.
 
Le premier soin des parens doit être de laisser à leurs enfans une éducation sainte, en travaillant de concert avec le pasteur pour les former à la vertu. Le premier soin des pasteurs doit être : 1. de veiller sur ces jeunes plantes; 2. de procurer à leurs paroisses des maîtres d’école vertueux, édifians, capables d’instruire les enfans; 3. d’établir des maîtresses d’école prudentes, pour élever les jeunes filles. Combien de filles et de mères vertueuses qui seraient perverties sans le secours d’une maîtresse d’école ! Les paroisses doivent avoir une singulière attention sur ce point. Ceux qui s’opposent à un établissement si important en répondront à Dieu. Si chacun doit contribuer à l’intérêt d’une communauté, le zèle de Dieu doit engager tout le monde à contribuer aussi à l’éducation de la jeunesse, surtout des jeunes filles.
 
\secundo Les personnes qui s’appliquent à l’instruction des enfans, doivent estimer cet emploi, dont plusieurs personnes, même de qualité, se sont fait honneur; et se souvenir qu’il n’y a aucun de ces enfans, quelque ignorant et pauvre qu’il soit, qui n’ait coûté à Jésus-Christ son sang. Ne vous rebutez point : les enfans ont plus de disposition à la vertu qu’on ne pense; les semences de piété qu’on jette dans ces jeunes cœurs, tôt ou tard produiront leurs fruits.
 
Si le Sauveur a témoigné de la tendresse pour les enfans, s’il les faisait approcher de sa personne, c’est parce que leur candeur et leur simplicité les rendaient plus dociles à sa parole, et plus propres au royaume des cieux\footnote{Sinite parvulos venire ad me, talium est enim regnum cœlorum. Matth. 19. 14.}.
 
Quoiqu’un enfant soit volage et dissipé, prenez garde, disait Jésus-Christ, de ne mépriser aucun de ces petits, parce que leurs anges voient la face de mon Père dans le ciel.
 
Nous devons nous-mêmes espérer de voir un jour dans le ciel cette face adorable, si nous avons du zèle pour les enfans, et si, à l’exemple des anges, nous sommes les gardiens de leur innocence.
 
Un saint homme disait à sa mort, qu’il lui semblait voir une nombreuse troupe d’enfans bienheureux qu’il avait instruits, qui lui venaient au-devant. Quelle joie à notre mort, si les âmes, au salut desquelles nous aurions contribué, venaient nous accompagner comme en triomphe jusque dans la gloire ! Ceux qui instruisent les autres dans la vertu, dit le Saint-Esprit, brilleront comme des astres dans l’éternité\footnote{Qui ad justitiam erudiunt multos, quasi stellæ in perpetuas æternitates. Dan. 12. 3.}.
 
\section{Résolutions}
 
\primo Ce sera surtout envers les enfans que j’aimerai à déployer mon zèle, parce qu’ils opposent d’ordinaire moins d’obstacles à la grâce, et que s’ils mènent une vie innocente, ils procurent à Dieu plus de gloire.
 
\secundo Mais en leur donnant des avis et des instructions, je regarderai soigneusement de rien leur dire qui puisse leur inspirer l’idée du mal, qu’ils ignorent peut-être encore.
 
\tertio Je fournirai aussi, autant qu’il dépendra de moi, pour leur éducation et pour leur faire apprendre des métiers convenables. 

\prayer{Mon Dieu, vous aimez les enfans; je désire être l'instrument de votre miséricorde, aidez-moi de votre grâce}
 
\end{document}
